\documentclass[aspectratio=169,11pt]{beamer}

\usepackage[T1]{fontenc}
\usepackage[utf8]{inputenc}
\usepackage[turkish,shorthands=off]{babel}
\usepackage{graphicx}
\usepackage{booktabs}
\usepackage{amsmath}
\usepackage{hyperref}

\usetheme{default}
\setbeamertemplate{navigation symbols}{}
\setbeamertemplate{footline}[frame number]

\definecolor{ustnavy}{HTML}{0B1F3A}
\definecolor{ustteal}{HTML}{1F6F8B}
\definecolor{ustaccent}{HTML}{8A2B0E}
\definecolor{ustgray}{HTML}{4A5560}
\definecolor{ustgreen}{HTML}{1B5E3B}

\setbeamercolor{structure}{fg=ustteal}
\setbeamercolor{frametitle}{fg=ustnavy}
\setbeamercolor{title}{fg=ustnavy}
\setbeamercolor{subtitle}{fg=ustgray}
\setbeamercolor{author}{fg=ustgray}
\setbeamercolor{institute}{fg=ustgray}
\setbeamercolor{item}{fg=ustteal}
\setbeamercolor{block title}{fg=white,bg=ustteal}
\setbeamercolor{block body}{fg=ustnavy,bg=ustnavy!4}
\setbeamercolor{block title alerted}{fg=white,bg=ustaccent}
\setbeamercolor{block body alerted}{fg=ustnavy,bg=ustaccent!8}
\setbeamercolor{block title example}{fg=white,bg=ustgreen}
\setbeamercolor{block body example}{fg=ustnavy,bg=ustgreen!8}
\setbeamerfont{title}{series=\bfseries,size=\Large}
\setbeamerfont{frametitle}{series=\bfseries,size=\large}

\newcommand{\shot}[2]{%
  \IfFileExists{#1}{%
    \includegraphics[width=#2\textwidth]{#1}%
  }{%
    \fbox{\begin{minipage}{0.88\textwidth}\centering\vspace{1.1cm}
    \textbf{Görsel bekleniyor}\\[0.35em]
    {\small\ttfamily\detokenize{#1}}\\[0.35em]
    {\small Lütfen dosyayı \texttt{user\_shots/} klasörüne ekleyiniz.}
    \vspace{1.1cm}\end{minipage}}%
  }%
}

\title{Universal Spectrum Toolkit}
\subtitle{MCA spektrum dosyalarının ROOT analiz nesnelerine dönüştürülmesi}
\author{Eren Yılmaz}
\institute{Teknik sunum}
\date{2026}

\begin{document}

%---------------------------------------------------------------------
\begin{frame}[plain]
  \titlepage
  \vspace{-0.4em}
  \begin{center}
    \small\textcolor{ustgray}{Sorun tanımı, yöntemsel yaklaşım, uygulama ve kullanım}
  \end{center}
\end{frame}

%---------------------------------------------------------------------
\begin{frame}{İçerik}
  \begin{enumerate}
    \item Problem tanımı: ROOT neden \texttt{.spe} dosyasını spektrum olarak işleyemez?
    \item Hatalı okuma: layout bilinmeden yapılan yorumun sonucu
    \item Yöntem: binary yapının dönüştürülmesi için kurulan mantık
    \item Matematiksel karar: layout hipotezlerinin skorlanması
    \item Uygulama çıktısı ve doğrulama
    \item Kullanım: menü tabanlı arayüz
  \end{enumerate}
\end{frame}

%---------------------------------------------------------------------
\begin{frame}{Problem tanımı}
  \begin{alertblock}{Temel sorun}
    ROOT bir analiz çerçevesidir. Beklediği veri modeli \texttt{.root},
    \texttt{TH1}, \texttt{TTree} ve yerleşik I/O tanımlarıdır.
    MCA üretimi \texttt{.spe}/\texttt{.chn} dosyaları ise ham bayt dizileridir.
  \end{alertblock}

  \vspace{0.5em}
  \textbf{Sonuç:}
  \begin{itemize}
    \item Dosyanın fiziksel olarak açılması, spektrumun doğru okunması anlamına gelmez.
    \item \texttt{.spe} tek bir uluslararası standart değildir.
    \item Header boyutu, sayısal tip ve bayt sırası üreticiye göre değişir.
  \end{itemize}
\end{frame}

%---------------------------------------------------------------------
\begin{frame}{Makro olmaksızın yapılan okuma denemesi}
  \begin{center}
    \shot{user_shots/03_blind_spe.png}{0.86}
  \end{center}
\end{frame}

%---------------------------------------------------------------------
\begin{frame}{Bu görüntü neyi kanıtlar?}
  \begin{columns}[T]
    \begin{column}{0.48\textwidth}
      \textbf{uint16 LE varsayımı}
      \begin{itemize}
        \item Yoğun gürültü zemini
        \item Fiziksel olarak anlamsız dikenler
        \item Spektrum yapısı oluşmaz
      \end{itemize}
    \end{column}
    \begin{column}{0.48\textwidth}
      \textbf{float32 LE varsayımı}
      \begin{itemize}
        \item $10^{18}$ mertebesinde değerler
        \item Görselleştirilemeyen çıktı
        \item Yanlış layout / yanlış hizalama
      \end{itemize}
    \end{column}
  \end{columns}

  \vspace{0.8em}
  \begin{block}{Çıkarım}
    Sorun ``dosyayı çizmek'' değildir. Sorun, bayt dizisine doğru
    \textbf{anlam haritasını} vermektir.
  \end{block}
\end{frame}

%---------------------------------------------------------------------
\begin{frame}{Yöntemsel fikir}
  \begin{exampleblock}{Merkezi yaklaşım}
    Dosyaya spektrum muamelesi yapmak yerine, önce binary düzeni
    sistematik hipotezlerle çözmek; doğrulanan düzeni koda kural olarak
    yerleştirmek.
  \end{exampleblock}

  \vspace{0.6em}
  \textbf{Üç adımlı düşünce:}
  \begin{enumerate}
    \item Ham dosyayı bayt düzeyinde incelemek
    \item Olası layout hipotezlerini üretmek ve skorlamak
    \item Kazanan layout'u kanal--counts dizisine, ardından ROOT nesnesine dönüştürmek
  \end{enumerate}
\end{frame}

%---------------------------------------------------------------------
\begin{frame}{Binary yapının çözülmesi}
  \begin{center}
    \shot{user_shots/04_binary_probe.png}{0.90}
  \end{center}
\end{frame}

%---------------------------------------------------------------------
\begin{frame}{Matematiksel karar modeli}
  Dosya boyutu $S$ verildiğinde her hipotez şu parametrelerle tanımlanır:

  \[
  H = S - N\cdot B
  \]

  \begin{itemize}
    \item $N$: varsayılan kanal sayısı (ör.\ $16384$)
    \item $B$: değer başına bayt (ör.\ \texttt{float32} için $B=4$)
    \item $H$: header ofseti (ör.\ $H=40$)
  \end{itemize}

  \vspace{0.4em}
  Payload, $[H,\ S)$ aralığında seçilen tip ve endianness ile çözülür.
  Elde edilen dizi fiziksel spektrum olma olasılığına göre skorlanır:
  \begin{itemize}
    \item sonlu değer oranı
    \item negatif olmama oranı
    \item düzgünlük / tutarlılık ölçütleri
  \end{itemize}
\end{frame}

%---------------------------------------------------------------------
\begin{frame}{Bu mantığın koda aktarılması}
  \begin{block}{Keşif aşaması}
    Probe çıktısı, örnek veri için en yüksek skorlu düzenin\\
    $H=40$, $N=16384$, \texttt{float32} olduğunu gösterir.
  \end{block}

  \vspace{0.5em}
  \begin{exampleblock}{Üretim kuralı}
    Bu düzen ``rastgele tahmin'' olarak bırakılmaz.\\
    Kodda doğrulanmış okuyucu kuralına dönüştürülür:\\
    header + kanal sayısı + sayısal tip + bayt sırası.
  \end{exampleblock}

  \vspace{0.5em}
  \textbf{Sonuç:} Yazılım artık dosyayı spekülasyonla değil,\\
  önceden doğrulanmış bir anlam modeli ile okur.
\end{frame}

%---------------------------------------------------------------------
\begin{frame}{Dönüşümün doğrulanması}
  \begin{center}
    \shot{user_shots/02_toolkit_match.png}{0.90}
  \end{center}
\end{frame}

%---------------------------------------------------------------------
\begin{frame}{Karar çıktısının anlamı}
  \begin{itemize}
    \item ASCII okuyucu eşleşmez: dosya metin değildir.
    \item GF3-benzeri okuyucu eşleşir: $40$ bayt header + \texttt{float32} LE.
    \item Kanal sayısı: $16384$
    \item Üretilen nesneler:
      \begin{itemize}
        \item \texttt{TH1D hSpectrumChannel}
        \item \texttt{TTree spectrum}
        \item \texttt{conversion\_metadata}
        \item ham bayt arşivi
      \end{itemize}
  \end{itemize}

  \vspace{0.5em}
  \begin{block}{Özet}
    Ham bayt dizisi, ROOT'un analiz edebileceği spektrum nesnesine
    kontrollü biçimde dönüştürülmüştür.
  \end{block}
\end{frame}

%---------------------------------------------------------------------
\begin{frame}{Başarılı dönüşüm çıktısı}
  \begin{columns}[c]
    \begin{column}{0.48\textwidth}
      \centering
      \shot{user_shots/05_converted_linear.png}{1.0}\\
      \scriptsize Lineer ölçek
    \end{column}
    \begin{column}{0.48\textwidth}
      \centering
      \shot{user_shots/06_converted_log.png}{1.0}\\
      \scriptsize Logaritmik ölçek
    \end{column}
  \end{columns}
\end{frame}

%---------------------------------------------------------------------
\begin{frame}{Gerçek veri doğrulaması: Co-60}
  \begin{columns}[c]
    \begin{column}{0.48\textwidth}
      \centering
      \shot{user_shots/08_co60_linear.png}{1.0}\\
      \scriptsize Lineer
    \end{column}
    \begin{column}{0.48\textwidth}
      \centering
      \shot{user_shots/07_co60_log.png}{1.0}\\
      \scriptsize Logaritmik
    \end{column}
  \end{columns}
\end{frame}

%---------------------------------------------------------------------
\begin{frame}[fragile]{Kullanım}
  Araç, ROOT komut bilgisi gerektirmeyen menü arayüzü ile kullanılır.

  \vspace{0.6em}
  \begin{verbatim}
cd UniversalSpectrumToolkit
./ust_menu.sh
  \end{verbatim}

  \vspace{0.4em}
  Menü seçenekleri:
  \begin{enumerate}
    \item Tek dosya dönüştürme
    \item Klasör genelinde toplu dönüştürme
    \item Örnek dosyalar ile deneme
    \item Yardım
    \item Çıkış
  \end{enumerate}
\end{frame}

%---------------------------------------------------------------------
\begin{frame}{Menü akışı}
  \begin{center}
    \shot{user_shots/01_menu.png}{0.88}
  \end{center}
\end{frame}

%---------------------------------------------------------------------
\begin{frame}{Menüdeki karar noktaları}
  \begin{itemize}
    \item Spektrumun ekranda çizilmesi: Evet / Hayır
    \item Ölçek seçimi: logaritmik Y / lineer Y
    \item Bilinmeyen formatta zorla generic deneme: varsayılan olarak kapalı
  \end{itemize}

  \vspace{0.7em}
  \begin{exampleblock}{Tasarım ilkesi}
    Kullanıcı yalnızca dosya yolunu ve görselleştirme tercihini belirtir.\\
    Format tanıma ve dönüşüm, doğrulanmış kurallar üzerinden yürütülür.
  \end{exampleblock}
\end{frame}

%---------------------------------------------------------------------
\begin{frame}{Sonuçlar}
  \begin{enumerate}
    \item ROOT, \texttt{.spe} içeriğini native spektrum modeli olarak bilmez.
    \item Layout bilinmeden yapılan okuma fiziksel olarak geçersiz sonuç üretir.
    \item Çözüm, binary düzenin hipotez--skor--doğrulama ile çözülmesidir.
    \item Doğrulanan düzen koda kural olarak aktarılır; çıktı ROOT analiz nesnesidir.
    \item Kullanım, menü arayüzü ile ROOT bilgisi olmadan gerçekleştirilebilir.
  \end{enumerate}
\end{frame}

%---------------------------------------------------------------------
\begin{frame}[plain]
  \begin{center}
    \vspace{1.8cm}
    {\Large\bfseries\textcolor{ustnavy}{Teşekkürler}}\\[1.0em]
    {\large\textcolor{ustteal}{Universal Spectrum Toolkit}}\\[1.2em]
    \textcolor{ustgray}{Sorularınız için hazırım.}
  \end{center}
\end{frame}

\end{document}
